\chapter*{Introduction}
\addcontentsline{toc}{chapter}{Introduction}
\markboth{Introduction}{Introduction}
\label{chap:introduction}

\noindent
Modern power grids are evolving faster than at any point in their history.
The integration of variable renewable energy sources such as wind, solar
photovoltaic, and run-of-river hydro has replaced the predictable output
of conventional thermal plants with generation profiles that fluctuate on
timescales of minutes to hours. At the same time, demand-side dynamics are
growing in complexity: electric vehicle charging, smart building management
systems, and responsive industrial loads all introduce new patterns of
uncertainty that deterministic point forecasts cannot adequately characterize.
In this environment, \emph{probabilistic forecasting}, which provides a
full predictive distribution of future electricity demand rather than a single
value, has transitioned from an academic research topic to an operational
necessity for grid operators, energy traders, and capacity planners alike.

\vspace{0.5cm}

\noindent
This project addresses the challenge of short-term probabilistic electricity
demand forecasting from a practical research perspective. Its starting point
is the recently published \emph{any-quantile} (AQ) framework of Smyl et
al.~\cite{smyl2025anyquantile}, which enables a single trained neural network
to generate forecasts at any arbitrary quantile level at inference time.
Applied to the N-BEATS architecture, this framework achieved a Continuous
Ranked Probability Score (CRPS) of 211.22 on the EMHIRES benchmark, 35
European countries, 48-hour horizon, 2018 test year, establishing a
compelling reference performance. The benchmark paper explicitly
identified exogenous meteorological features and more sophisticated
quantile-sampling strategies as promising directions for improvement.

\vspace{0.5cm}

\noindent
The present report documents our systematic investigation of these and
other enhancement strategies. Through a carefully controlled ablation study
spanning more than twenty configurations, we demonstrate that integrating
weather and calendar exogenous features into AQ-NBEATS reduces CRPS by
approximately 20\%, surpassing even the stronger AQ-ESRNN variant from
the same paper. A curriculum learning approach combining beta-distributed
tail sampling with adaptive quantile scheduling achieves our best
multi-seed result of CRPS = 171.3. Finally, post-hoc Conformalized
Quantile Regression (CQR) corrects residual calibration gaps, achieving
nominal 95\% empirical coverage even under seasonal distribution shift.

\vspace{0.5cm}

\noindent
\textbf{Structure of the report.}
Chapter~1 reviews the technical background and related literature.
Chapter~2 analyses and specifies the system requirements.
Chapter~3 presents the software and architectural design.
Chapter~4 details the experimental methodology.
Chapter~5 presents implementation details and empirical results with
comprehensive analysis. A final chapter concludes and outlines future
research perspectives.
